\begin{lemma}[The truncated kernel integral varies little over a small ball]
Fix $n \geq 1$ and $j \in \{1,\ldots,n\}$, let $\lambda$ be Lebesgue measure on
$\mathbb{R}^n$, let $\gamma_n$ be the Gaussian measure of
\noderef{GaussMeasure} and let $K_{\gamma,j}$ be the kernel of
\noderef{GaussRieszKernel}.  Let $R \geq 1$.  Then there is a constant
$C \in (0,\infty)$, depending only on $n$, $j$ and $R$, such that the following
holds.

Let $g : \mathbb{R}^n \to \mathbb{R}$, let $\Phi$ and $N$ be real, let
$x,z \in \mathbb{R}^n$ and let $\varepsilon$ be real, and suppose
\[
  |x| \leq R, \qquad 0 < \varepsilon \leq 1, \qquad |x-z| < \tfrac12\varepsilon ,
  \qquad 0 \leq \Phi ,
\]
and that
\begin{enumerate}
\item[(a)] $g$ is $\gamma_n$-integrable;
\item[(b)] $g = 0$ $\lambda$-almost everywhere on $\{y : |x-y| \leq \varepsilon\}$;
\item[(c)] $\displaystyle\int_{\overline B(x,\rho)}^{*}|g|\,d\lambda
  \leq \Phi\cdot\lambda\bigl(\overline B(x,\rho)\bigr)$ for every real
  $\rho \in (0,2]$, the left side being the lower Lebesgue integral in
  $[0,+\infty]$ and the right side read there too; and
\item[(d)] $y \mapsto e^{-|y|^2}|g(y)|$ is Lebesgue integrable on
  $\mathbb{R}^n$, with
  $\displaystyle\int_{\mathbb{R}^n}e^{-|y|^2}|g(y)|\,d\lambda(y) \leq N$.
\end{enumerate}
Then
\[
  \Bigl|\int_{\mathbb{R}^n}K_{\gamma,j}(x,y)g(y)\,dy
      - \int_{\mathbb{R}^n}K_{\gamma,j}(z,y)g(y)\,dy\Bigr|
    \;\leq\; C\,(\Phi + N) ,
\]
both integrals being absolutely convergent Bochner integrals.

\emph{No hypothetical operator is needed here, and that is not an oversight.}
A reader coming from a Cotlar-type argument will expect the two whole-space
integrals to be identified with coordinates of a transform, and hence expect an
abstract $T$ with its two properties to be carried along.  They are not needed:
the only thing required of the two integrals is that they converge absolutely,
and that comes from clause (1)(b) of \noderef{GaussRieszKernelFarMajorant} at
the pairs $(R,\varepsilon)$ and $(R+1,\varepsilon/2)$, whose constants are
irrelevant here because only the integrability half of that clause is used.  A
consumer that also wants to identify these integrals with $(Tg)_j$ must invoke
\noderef{GaussRieszOffDiagonalL1} separately; that identification is not part of
this lemma.

\emph{Why (c) is a lower Lebesgue integral, and why (d) is a hypothesis.}
Hypothesis (c) is stated with $\int^{*}$ because that is the shape a consumer
has available: it is what the definition of the local maximal function of
\noderef{LocalMaximalFunction} as a supremum over $\rho \in (0,2]$ supplies,
after multiplying by the finite positive number
$\lambda(\overline B(x,\rho))$.  Hypothesis (d) is assumed rather than derived
from (a) because the Gaussian-to-Lebesgue density conversion that supplies it
--- $\int e^{-|y|^2}|g|\,d\lambda = \pi^{n/2}\int|g|\,d\gamma_n$ --- is used by
the intended consumer in a second, independent place as well, so it belongs in
the consumer once rather than being re-derived here.

\emph{The constant depends on $R$, necessarily.}  $C$ may be taken to be
$\max(C_d,\,C_a+C_b)$ with $C_d$ the constant of
\noderef{GaussRieszKernelHormanderDyadicBound} at $R$ and $C_a$, $C_b$ the
constants of clause (1) of \noderef{GaussRieszKernelFarMajorant} at the pairs
$(R,1)$ and $(R+1,\tfrac12)$; all three degrade in $R$, and no uniformity in
$n$, $j$ or $R$ is claimed.  Note that $C$ does not depend on $\varepsilon$:
that is the point of the lemma, and it is why the two far-field constants are
taken at the fixed radii $1$ and $\tfrac12$ rather than at $\varepsilon$ and
$\varepsilon/2$.
\end{lemma}
\begin{proof}
Let $C_d > 0$ be the constant supplied by
\noderef{GaussRieszKernelHormanderDyadicBound} at $R$, and let $C_a, C_b > 0$ be
the constants supplied by clause (1) of \noderef{GaussRieszKernelFarMajorant} at
the pairs $(R,1)$ and $(R+1,\tfrac12)$ respectively; the hypotheses $R > 0$,
$1 > 0$, $R+1 > 0$ and $\tfrac12 > 0$ all hold because $R \geq 1$.  Set
$C = \max(C_d,\,C_a+C_b) \in (0,\infty)$, which depends only on $n$, $j$ and
$R$.  Fix $g$, $\Phi$, $N$, $x$, $z$, $\varepsilon$ and the hypotheses.

\emph{Step 0: two preliminaries.}  Since $e^{-|y|^2}|g(y)| \geq 0$ for every
$y$, its integral is nonnegative, so (d) gives $N \geq 0$.  And
$|z| - |x| \leq |z-x| = |x-z| < \varepsilon/2 \leq \tfrac12$, so
\begin{equation*}
  |z| \leq R + 1 . \tag{1}
\end{equation*}

\emph{Step 1: three regions.}  Put
\[
  S = \{y : |x-y| > \varepsilon\}, \qquad
  A_1 = \{y : \varepsilon < |x-y| \leq 1\}, \qquad
  A_2 = \{y : |x-y| > 1\} .
\]
All three are Borel, being preimages of an interval under the continuous map
$y \mapsto |x-y|$.  They satisfy $S = A_1 \cup A_2$ with $A_1 \cap A_2 =
\emptyset$: the union is contained in $S$ because $\varepsilon \leq 1$, the
reverse inclusion holds by splitting on whether $|x-y| \leq 1$, and the
intersection is empty because $y \in A_1$ gives $|x-y| \leq 1$ while $y \in A_2$
gives $|x-y| > 1$.

\emph{Step 2: both kernel products are integrable on $S$.}  By (a), $g$
satisfies the hypothesis of clause (1)(b) of
\noderef{GaussRieszKernelFarMajorant}.  Applying that clause at the pair
$(R,\varepsilon)$ --- legitimate since $\varepsilon > 0$ --- to $g$ and to the
point $x$, which has $|x| \leq R$, gives that
$y \mapsto K_{\gamma,j}(x,y)g(y)$ is Lebesgue integrable on
$\{y : |x-y| > \varepsilon\} = S$.  Applying the same clause at the pair
$(R+1,\varepsilon/2)$ to $g$ and to the point $z$, which has $|z| \leq R+1$ by
$(1)$, gives that $y \mapsto K_{\gamma,j}(z,y)g(y)$ is Lebesgue integrable on
$\{y : |z-y| > \varepsilon/2\}$.  Now
\begin{equation*}
  S \subseteq \{y : |z-y| > \varepsilon/2\} , \tag{2}
\end{equation*}
because for $y \in S$, writing $x-y = (x-z)+(z-y)$ and using the triangle
inequality, $|z-y| \geq |x-y| - |x-z| > \varepsilon - \varepsilon/2 =
\varepsilon/2$.  A function integrable on a set is integrable on every
measurable subset of it, so $y \mapsto K_{\gamma,j}(z,y)g(y)$ is Lebesgue
integrable on $S$ as well, and both products are integrable on each of
$A_1 \subseteq S$ and $A_2 \subseteq S$.

\emph{Step 3: the whole-space integrals collapse onto $S$.}  Let $w$ be either
$x$ or $z$.  The complement of $S$ is exactly $\{y : |x-y| \leq \varepsilon\}$,
on which $g$ vanishes $\lambda$-almost everywhere by (b); hence
$K_{\gamma,j}(w,y)g(y) = 0$ for $\lambda$-almost every $y \notin S$.  A
function integrable on a measurable set and vanishing almost everywhere off it
is integrable on the whole space with the same integral, so by Step 2
\begin{equation*}
  \int_{\mathbb{R}^n}K_{\gamma,j}(w,y)g(y)\,dy
    = \int_S K_{\gamma,j}(w,y)g(y)\,dy . \tag{3}
\end{equation*}
In particular both whole-space integrals in the conclusion are absolutely
convergent.

\emph{Step 4: the splitting.}  By Step 1, Step 2 and additivity of the integral
over a disjoint union of measurable sets,
$\int_S = \int_{A_1} + \int_{A_2}$ for each of the two products.  Subtracting
and using linearity of the integral on $A_1$ --- legitimate because both
products are integrable there --- we get from $(3)$
\begin{equation*}
  \int K_{\gamma,j}(x,y)g\,dy - \int K_{\gamma,j}(z,y)g\,dy
    = \int_{A_1}\bigl(K_{\gamma,j}(x,y)-K_{\gamma,j}(z,y)\bigr)g(y)\,dy
      + \Bigl(\int_{A_2}K_{\gamma,j}(x,y)g\,dy - \int_{A_2}K_{\gamma,j}(z,y)g\,dy\Bigr) .
  \tag{4}
\end{equation*}

\emph{Step 5: the near annulus.}  The hypotheses of
\noderef{GaussRieszKernelHormanderDyadicBound} at $R$ are met with this $g$,
with $\Phi$ read in $[0,+\infty]$, and with these $x$, $z$, $\varepsilon$: that
node requires $|x| \leq R$, $0 < \varepsilon \leq 1$,
$|x-z| < \varepsilon/2$ and exactly hypothesis (c).  It therefore gives
\[
  \int_{A_1}^{*}\bigl|K_{\gamma,j}(x,y)-K_{\gamma,j}(z,y)\bigr|\,|g(y)|\,d\lambda(y)
    \leq C_d\,\Phi
\]
in $[0,+\infty]$.  For any function $h$ and any measurable set $E$ one has
$\bigl|\int_E h\,d\lambda\bigr| \leq \int_E^{*}|h|\,d\lambda$, the right side
read in $[0,+\infty]$ and the inequality understood after passing the right side
to $[0,\infty)$ when it is finite; here it is finite, being at most $C_d\Phi$
with $C_d, \Phi$ real and $\Phi \geq 0$.  Applying this with $E = A_1$ and
$h(y) = (K_{\gamma,j}(x,y)-K_{\gamma,j}(z,y))g(y)$, whose absolute value is
$|K_{\gamma,j}(x,y)-K_{\gamma,j}(z,y)|\,|g(y)|$,
\begin{equation*}
  \Bigl|\int_{A_1}\bigl(K_{\gamma,j}(x,y)-K_{\gamma,j}(z,y)\bigr)g(y)\,dy\Bigr|
    \leq C_d\,\Phi . \tag{5}
\end{equation*}

\emph{Step 6: the far annulus.}  Let $y \in A_2$, so $|x-y| > 1$ and in
particular $1 \leq |x-y|$.  Since $|x| \leq R$, clause (1)(a) of
\noderef{GaussRieszKernelFarMajorant} at the pair $(R,1)$ gives
$|K_{\gamma,j}(x,y)| \leq C_ae^{-|y|^2}$, hence
$|K_{\gamma,j}(x,y)g(y)| \leq C_a\,e^{-|y|^2}|g(y)|$.  Both of these are
integrable on $A_2$ --- the left by Step 2 and the right by (d) --- so,
integrating the pointwise bound and then using that $e^{-|y|^2}|g(y)| \geq 0$
is integrable on $\mathbb{R}^n$, so that its integral over the subset $A_2$ is
at most its integral over $\mathbb{R}^n$,
\begin{equation*}
  \Bigl|\int_{A_2}K_{\gamma,j}(x,y)g(y)\,dy\Bigr|
    \leq \int_{A_2}\bigl|K_{\gamma,j}(x,y)g(y)\bigr|\,d\lambda
    \leq C_a\int_{\mathbb{R}^n}e^{-|y|^2}|g|\,d\lambda
    \leq C_a\,N , \tag{6}
\end{equation*}
the last step by (d).  Symmetrically, for $y \in A_2$ we have $|z| \leq R+1$ by
$(1)$ and
\[
  |z-y| \geq |x-y| - |x-z| > 1 - \tfrac12\varepsilon \geq 1 - \tfrac12
    = \tfrac12 ,
\]
using $\varepsilon \leq 1$; so clause (1)(a) at the pair $(R+1,\tfrac12)$ gives
$|K_{\gamma,j}(z,y)| \leq C_be^{-|y|^2}$, and the same three steps give
\begin{equation*}
  \Bigl|\int_{A_2}K_{\gamma,j}(z,y)g(y)\,dy\Bigr| \leq C_b\,N . \tag{7}
\end{equation*}

\emph{Step 7: conclusion.}  Taking absolute values in $(4)$ and using the
triangle inequality twice, then $(5)$, $(6)$ and $(7)$,
\[
  \Bigl|\int K_{\gamma,j}(x,y)g\,dy - \int K_{\gamma,j}(z,y)g\,dy\Bigr|
    \leq C_d\Phi + \bigl(C_aN + C_bN\bigr) .
\]
Finally $C_d \leq C$ and $C_a + C_b \leq C$ by definition of $C$, and
$\Phi \geq 0$ and $N \geq 0$, so
$C_d\Phi \leq C\Phi$ and $(C_a+C_b)N \leq CN$, whence the right-hand side is at
most $C(\Phi+N)$.
\end{proof}
